% TOP_PROBLEM: {"id": 30000209, "title": "Atiyah Conjecture on $L^2$-Betti Numbers", "category_id": 7, "category": {"id": 7, "name": "topology", "display_name": "Topology", "description": "Properties preserved under continuous deformations.", "slug": "topology", "order_index": 7, "created_at": "2026-07-31T15:26:25.671Z"}, "status": "partially_solved"}
% FIRST_POSTED: 2026-09-27
% !TeX program = lualatex
% ATTEMPT_STATUS: unresolved
\documentclass[11pt]{article}
\usepackage[margin=25mm]{geometry}
\usepackage{fontspec}
\setmainfont{Latin Modern Roman}
\usepackage{amsmath,amssymb,mathtools}
\usepackage{unicode-math}
\setmathfont{Latin Modern Math}
\usepackage{xurl,hyperref}
\hypersetup{colorlinks=true,urlcolor=blue}
\setcounter{secnumdepth}{0}
\setlength{\parindent}{0pt}
\setlength{\parskip}{5pt}
\setlength{\emergencystretch}{3em}
\begin{document}
\section*{\texorpdfstring{Atiyah Conjecture on $L^2$-Betti Numbers}{Atiyah Conjecture on \$L\textasciicircum{}2\$-Betti Numbers}}\label{30000209}
\subsection{Definitions and mathematical statement}
Let $X$ be a finite CW-complex and $G=\pi_1(X)$ torsion-free. Lifting cells to its universal cover gives a Hilbert chain complex with spaces $C_i^{(2)}(\widetilde X)\cong\ell^2(G)^{c_i}$ and bounded cellular boundary maps $d_i$. Its reduced $L^2$ homology is
\[
 \mathcal H_i^{(2)}=\ker d_i\big/\overline{\operatorname{im}d_{i+1}},
 \qquad b_i^{(2)}(\widetilde X)=\dim_{\mathcal N(G)}\mathcal H_i^{(2)}.
\]
The von Neumann dimension is computed using the canonical group trace on the orthogonal projection representing this Hilbert module, not an ordinary vector-space dimension. The selected Atiyah assertion is $b_i^{(2)}(\widetilde X)\in{\mathbb{Z}}_{\ge0}$ for every $i$ and every such $X$. No assertion about groups with torsion or denominator bounds for their invariants is substituted.
\par\smallskip\textit{Formulation and concept references:} \href{https://arxiv.org/abs/math/0001101}{[S1]}.

\subsection{Short English statement}
For every finite cell complex with torsion-free fundamental group, must all von Neumann dimensions of its square-integrable homology be integers?
\subsection{Research attempt: Fourier reduction and the integrality obstruction}
\textbf{Status and coefficient scope.} The general torsion-free assertion is not established here. The inherited source [S1] studies integrality through classes of groups. A current primary manuscript of Fisher and Ng [E1], submitted in June 2026, treats specified outer automorphism groups and torsion-free finite-index subgroups, with stronger algebraic conclusions in that setting. It is not a proof for every torsion-free group. The work below gives complete calculations for free abelian fundamental groups and graphs, and distinguishes ordinary Euler characteristic from the individual von Neumann dimensions.

\textbf{The dimension bookkeeping.} For a bounded $G$-equivariant map
\[
 A:\ell^2(G)^a\longrightarrow\ell^2(G)^b,
\]
put $r(A)=\dim_{\mathcal N(G)}\overline{\operatorname{im}A}$. The basic Hilbert-module rank identity is
\[
 \dim_{\mathcal N(G)}\ker A+r(A)=a.
 \tag{334.1}
\]
One way to justify it is to use polar decomposition. The partial isometry in $A=U|A|$ identifies the orthogonal complement of $\ker A$ with the closure of the image; its initial and final projections have the same trace. The trace of the identity on $\ell^2(G)^a$ is $a$.

Because $d_id_{i+1}=0$, the closed image of $d_{i+1}$ lies in the closed kernel of $d_i$. Subtracting their dimensions gives
\[
 b_i^{(2)}=c_i-r(d_i)-r(d_{i+1}).
 \tag{334.2}
\]
In particular the alternating sum telescopes:
\[
 \sum_i(-1)^i b_i^{(2)}=\sum_i(-1)^i c_i=\chi(X).
 \tag{334.3}
\]
Equation (334.3) is valid without torsion-freeness. It cannot prove that each summand is an integer: cancellation of nonintegral numbers in an alternating sum is entirely possible. No assertion that an arbitrary numerical example is realized by a chain complex is needed to see this logical limitation.

\textbf{Lemma 334.1 (Laurent polynomial zero sets).} A nonzero complex Laurent polynomial on the torus $\mathbb T^d$ has zero set of Haar measure zero.

\emph{Proof.} Multiplying by a monomial reduces to an ordinary polynomial without changing its zero set on the torus. For $d=1$, a nonzero polynomial has finitely many roots. Induct on $d$, writing
\[
 f(z_1,\ldots,z_d)=\sum_{j=0}^m a_j(z_1,\ldots,z_{d-1})z_d^j.
\]
Choose a coefficient $a_j$ which is not the zero polynomial. By induction the set where this coefficient vanishes has measure zero. Outside that exceptional set, the polynomial in $z_d$ is nonzero and has finitely many zeros. Fubini's theorem now gives measure zero for the zero set of $f$. $\square$

\textbf{Theorem 334.2 (the free abelian case).} If $\pi_1(X)\cong\mathbb Z^d$, all the reduced $L^2$-Betti numbers in the selected statement are integers.

\emph{Proof.} For $d=0$ the group is trivial, the chain complex is finite-dimensional over $\mathbb C$, and there is nothing to prove. For $d\geq1$, Fourier transform identifies $\ell^2(\mathbb Z^d)$ with $L^2(\mathbb T^d)$. A matrix over the group ring becomes pointwise multiplication by a matrix $A(z)$ of Laurent polynomials. Let $r$ be its rank over the rational function field $\mathbb C(z_1,\ldots,z_d)$. Every $(r+1)$-minor is identically zero, and, if $r>0$, at least one $r$-minor is a nonzero Laurent polynomial. Lemma 334.1 shows
\[
 \operatorname{rank}_{\mathbb C}A(z)=r
 \quad\text{for almost every }z\in\mathbb T^d.
\]
The kernel projection is pointwise the finite-dimensional orthogonal projection onto $\ker A(z)$. This can also be obtained as the pointwise limit of $(I+nA(z)^*A(z))^{-1}$, making measurability explicit. The group trace is the integral of the ordinary matrix trace. Hence
\[
 \dim_{\mathcal N(\mathbb Z^d)}\ker A
    =\int_{\mathbb T^d}\dim_{\mathbb C}\ker A(z)\,dz
    =a-r.
 \tag{334.4}
\]
Thus each boundary map has integer von Neumann rank. Substitute these ranks into (334.2). Nonnegativity follows because the image being subtracted is a submodule of the kernel. $\square$

The proof is about almost-everywhere rank, not rank at a specially selected character of $\mathbb Z^d$. A boundary matrix may drop rank at $z=(1,\ldots,1)$ without affecting the trace integral.

\textbf{Two explicit complexes.} For the standard $d$-torus, Fourier-transformed cellular chains form the Koszul complex on $z_1-1,\ldots,z_d-1$. At a point where $z_j\ne1$, wedging by the $j$th basis vector and dividing by $z_j-1$ gives a chain contraction, with the usual exterior-algebra signs. The identity follows from the anticommutation relation between contraction and exterior multiplication. Every point except $(1,\ldots,1)$ has such a coordinate. Therefore the complex is exact almost everywhere, and all its reduced $L^2$-Betti numbers vanish for $d\geq1$.

This fiberwise contraction is not asserted to be a bounded Hilbert-space operator: the denominators grow near $z_j=1$. The almost-everywhere rank calculation, which is sufficient for (334.4), avoids that false boundedness claim. It also explains why closure of the image in the definition is essential.

For a contrasting example take
\[
 X=S^1\vee\bigvee_{j=1}^m S^2.
\]
Its fundamental group is $\mathbb Z$. The lifted chain complex has $c_0=c_1=1$, $c_2=m$, with $d_1=t-1$ and $d_2=0$. Multiplication by $z-1$ has zero kernel and dense image in $L^2(\mathbb T)$, because it is nonzero almost everywhere. Consequently
\[
 b_0^{(2)}=b_1^{(2)}=0,\qquad b_2^{(2)}=m.
\]
Thus torsion-free fundamental group does not mean that all positive-degree $L^2$-homology vanishes. The cell complex, as well as the group, enters the selected statement.

\textbf{A nonabelian calculation for graphs.} Let $X$ be a bouquet of $r\geq2$ circles, with $G=F_r$. There is one vertex orbit and $r$ edge orbits. The orthogonal complement of the closure of $\operatorname{im}d_1$ is $\ker d_1^*$. An element of this kernel is a function on $G$ invariant under all free generators and hence under all left translations. It is constant, and a constant square-summable function on an infinite group is zero. Therefore $r(d_1)=1$. There are no two-cells, and (334.2) gives
\[
 b_0^{(2)}=0,\qquad b_1^{(2)}=r-1.
\]
The universal cover is a tree and has zero ordinary first homology. Its positive reduced $L^2$-Betti number is not a contradiction: the chain spaces and admissible infinite square-summable chains differ from the finite cellular chains of ordinary homology.

\textbf{Why the torsion hypothesis cannot simply be removed.} For a finite group $G$ of order $q$, the projection onto the constant functions in $\ell^2(G)$ is convolution by
\[
 e=\frac1q\sum_{g\in G}g.
\]
Its group trace is the coefficient of the identity, namely $1/q$. A connected finite complex with this fundamental group has $b_0^{(2)}=1/q$. For $G=C_p$, a presentation complex supplies such an example. It lies outside the torsion-free hypothesis and therefore is not a counterexample to the selected assertion.

This normalization also explains finite-cover diagnostics. For the $N$-fold cyclic cover of $S^1\vee\bigvee_m S^2$, the ordinary Betti numbers are $1,1,mN$. Dividing by $N$ gives $1/N,1/N,m$, approaching the preceding $L^2$ values. Fractional values in a finite approximation do not themselves refute integrality of the limit.

\textbf{An algebraic route and its unproved premise.} A sufficient strengthening would say that the kernel dimension of every finite matrix over $\mathbb ZG$ is integral. Equations (334.1)--(334.2) immediately imply the selected statement from that matrix assertion. The reverse reduction, involving realization and coefficient conventions, is not needed here and is not silently assumed.

The Fourier proof works because determinants and minors over a commutative Laurent polynomial ring control rank outside a null set. For a general torsion-free group, its group ring need not be commutative, and there is no available torus of scalar characters on which all matrix ranks can be read off in this way. Replacing the group ring by a division ring would require proving an appropriate embedding and compatibility with von Neumann rank; that is major additional structure, not a consequence of torsion-freeness established here.

\textbf{Checks and outcome.} Exact finite diagnostics compute cyclic-cover boundary ranks and the normalized Betti numbers just described. They also evaluate selected Laurent minors at rational points to check the explicit examples; the almost-everywhere theorem is proved above, not inferred from samples. The attempt establishes integrality for free abelian fundamental groups and computes the graph and wedge examples. The unresolved step is a general integer-rank theorem for the noncommutative cellular boundary operators associated with arbitrary torsion-free groups. The selected assertion remains unresolved by this work.

\subsection{Sources}
\begin{itemize}
\item[S1]T. Schick, Math. Ann. 317 (2000), 727-750.
\url{https://arxiv.org/abs/math/0001101}.
\textit{Formulation source.}
\item[E1]Sam P. Fisher and Andrew Ng, \emph{Outer automorphism groups and the Atiyah Conjecture}, arXiv:2606.19606, submitted 17 June 2026.
\url{https://arxiv.org/abs/2606.19606}.
\textit{Primary source consulted 27 September 2026; the specified outer automorphism groups and finite-index hypotheses are retained.}
\end{itemize}
\end{document}
